\begin{center}
\textbf{\large{Abstract}}
\end{center}

\addcontentsline{toc}{chapter}{Abstract}
\begin{small}

This report presents the analysis, design, and implementation of \textbf{Fitty}, an Android mobile application that supports personalized workout planning and health tracking. The main problem addressed in this work is the development of a system that goes beyond a simple exercise catalog by covering the full user journey, including authentication, onboarding, starter plan generation, exercise exploration, workout execution, progress tracking, and user engagement through notifications.

To address this problem, the system is implemented using a layered architecture. Kotlin and Jetpack Compose are used for the presentation layer, use cases and repositories are used for business logic, Room and DataStore are employed for local persistence, and Firebase provides cloud-based services for authentication, data storage, media storage, and push notifications. A notable technical aspect of the system is the adoption of an \textit{offline-first} approach for the exercise module, where exercise metadata is synchronized from Firestore, validated, normalized, and stored in Room so that the user interface can access it more efficiently and reliably. In addition, the application supports dynamic content management through Firebase Content, allowing onboarding options, practice categories, starter plan templates, and behavioral configurations to be updated without modifying the UI source code.

The outcome of the report is a functional prototype of a personalized fitness support application. The implemented system includes the main modules for sign-up and sign-in, guest access, onboarding and plan preview, workout planning, exercise browsing by muscle category, exercise detail and media viewing, quick workouts and scheduled workouts, basic progress tracking, user profile management, local reminders, and Firebase Cloud Messaging notifications. The project also provides unit tests for several important business components such as starter plan generation, exercise synchronization policy, workout completion, and progress tracking.

Overall, the report demonstrates that combining a modern Android architecture with cloud services and local persistence is a suitable approach for building personalized health and fitness applications. The achieved results provide a solid foundation for future extensions such as more advanced recommendation logic, richer progress visualization, and a more complete real-world user experience.

\vspace*{1cm}
\textbf{Keywords}: Android, Jetpack Compose, Firebase, offline-first, personalized fitness application
\end{small}
